\documentclass[a4paper,11pt]{article}
\usepackage{amsmath,amssymb,amsthm, tikz,titlesec,hyperref,esint,braket, graphicx, mathrsfs, enumitem}
\usepackage[a4paper,margin=2cm]{geometry}
\linespread{1.3}
\newtheorem{claim}{Claim}[section]

\newcommand\name{Park Chanwoo}   % Name of the student
\newcommand\university{KAIST} % Name of the university
\newcommand\department{Physics} % Name of the department
\newcommand\studentid{20230297} % Student ID
\newcommand\s{\,\;}


\newenvironment{solution}[1]
  {\renewcommand\qedsymbol{$\square$}\begin{proof}[\textbf{Solution#1}]}
  {\end{proof}}
\newenvironment{note}
  {\renewcommand\qedsymbol{$\blacksquare$}\begin{proof}[\textnormal{\textbf{note}}]}
  {\end{proof}}

\title{KAIST\\2026 Special Topics in Physics quantum Electronic Devices I: Interferometers and Fractional Particles\\
Homework 6\bigskip}
\author{\textbf{\Large \name} \\
% University: \university\\
Department: \department\\
Student ID: \studentid}
\date{\today}

\begin{document}
\thispagestyle{empty}
\maketitle
\tableofcontents
\titleformat{\section}[frame]{\pagebreak}{\filright
\footnotesize  \enspace \textsf{KAIST --- Interferometers and Fractional Particles 2026 Spring}\enspace}{6pt}{\Large\bfseries\filcenter}
\setcounter{secnumdepth}{1}
\setcounter{tocdepth}{2}

\newcommand{\tr}{\operatorname{tr}}
\newcommand{\sgn}{\operatorname{sgn}}
\newcommand{\supp}{\operatorname{supp}}
\renewcommand{\Re}{\operatorname{Re}}
\renewcommand{\Im}{\operatorname{Im}}
\newcommand{\boltz}{k_{\mathrm{B}}}

\section{Kac Moody commutation relation}

Start from
\begin{equation}
  n_{R,q} = \sum_k :c^\dagger_{R,k+q}c_{R,k}: \qquad n_{R,-q} = \sum_k :c^\dagger_{R,k-q}c_{R,k}:
\end{equation}

since the contracted term does not affect to the commutator ($[A + c, B] = [A, B] + [c, B] = [A, B]$ for any complex number $c$), the commutator can be calculated by:
\begin{equation}
  [n_{R,-q}, n_{R,q}] = \sum_{k, k'} [c^\dagger_{R,k-q}c_{R,k}, c^\dagger_{R,k'+q}c_{R,k'}].
\end{equation}
Apply the identity
\begin{equation}
  [AB, CD] = A\{B, C\}D  - AC\{B, D\} + \{A, C\}DB - C\{A, D\}B
\end{equation}
with $A = c^\dagger_{R,k-q}, B = c_{R,k}, C = c^\dagger_{R,k'+q}, D = c_{R,k'}$.

The anticommutators are:
\begin{gather}
  \{B, C\} = \{c_{R,k}, c^\dagger_{R,k'+q}\} = \delta_{k, k'+q},\\
  \{B, D\} = \{c_{R,k}, c_{R,k'}\} = 0,\\
  \{A, C\} = \{c^\dagger_{R,k-q}, c^\dagger_{R,k'+q}\} = 0,\\
  \{A, D\} = \{c^\dagger_{R,k-q}, c_{R,k'}\} = \delta_{k-q, k'}.
\end{gather}

Gathering the terms survive,
\begin{align}
  [n_{R,-q}, n_{R,q}] &= \sum_{k, k'} [\delta_{k, k'+q} c^\dagger_{R,k-q} c_{R,k'} -\delta_{k-q, k'} c^\dagger_{R,k'+q} c_{R,k}]\\
  &= \sum_{k} [c^\dagger_{R,k-q} c_{R,k-q} - c^\dagger_{R,k} c_{R,k}]
\end{align}

To resolve $\infty-\infty$ problem, separate the operators into normal ordered part and ground state expectation value:
\begin{equation}
  [n_{R,-q}, n_{R,q}] = \sum_{k} [:c^\dagger_{R,k-q} c_{R,k-q}: - :c^\dagger_{R,k} c_{R,k}:] + \sum_k [\langle c^\dagger_{R,k-q} c_{R,k-q}\rangle_0 - \langle c^\dagger_{R,k} c_{R,k}\rangle_0].
\end{equation}

Since the first term converges absolutely, shift the terms and compensate term-by-term. After that,
\begin{equation}
  [n_{R,-q}, n_{R,q}] = \sum_k [\langle c^\dagger_{R,k-q} c_{R,k-q}\rangle_0 - \langle c^\dagger_{R,k} c_{R,k}\rangle_0] = \sum_k (n_{k-q}^0 - n_k^0)
\end{equation}

Since the ground state of non-interacting fermions is fermi sea, $n_k^0 = \Theta(k_F - k)$. For large or small enough $k$, the term $n_{k-q}^0 - n_k^0$ is zero because both $n_{k-q}^0$ and $n_k^0$ are zero or one respectively. For simplicity, assume $q > 0$. The only non-zero contribution comes from the region near Fermi points, where $n_{k-q}^0 - n_k^0$ is one for $k -q  < k_F$ and $k > k_F$. The sum is given by the number of $k \in \frac{L}{2\pi}\mathbb{Z}/L\mathbb{Z}$. Since the length of the region $k_F < k < k_F + q$ is $q$, the number of $k$ in the region is $\frac{L}{2\pi}q$. Hence, we get
\begin{equation}
  [n_{R,-q}, n_{R,q}] = \frac{qL}{2\pi}.
\end{equation}

For $q < 0$, the region with non-zero contribution is $k_F + q < k < k_F$, same length as the previous case, so we get the same result.


\section{1D periodic tight-binding chain}

\subsection{(2a)}

Consider 1D periodic tight-binding chain with $L$ sites. The Hamiltonian is given by
\begin{eqnarray}
  H = \epsilon_0\sum_{i=0}^{L} c_i^\dagger c_i - t \sum_{i=0}^{L} (c_i^\dagger c_{i+1} + c_{i+1}^\dagger c_i),
\end{eqnarray}
the periodic boundary condition is imposed, i.e., $c_{L+1} = c_1$. 

Since the system has lattice translation symmetry, it is useful to perform Fourier transformation:
\begin{eqnarray}
  c_k^\dagger = \frac{1}{\sqrt{L}} \sum_{j=0}^{L} e^{ikj} \tilde{c}_j^\dagger, \quad k \in \frac{2\pi}{L}\mathbb{Z}/L\mathbb{Z}
\end{eqnarray}

Explicitly, the transformation is given by
\begin{align}
  \sum_{i=0}^{L} c_i^\dagger c_i &= \sum_{i=0}^{L} \sum_{k, k'} \frac{1}{L} e^{i(k-k')i} \tilde{c}_k^\dagger \tilde{c}_{k'} = \sum_{k} \tilde{c}_k^\dagger \tilde{c}_k,\\
  \sum_{i=0}^{L} c_i^\dagger c_{i+1} &= \sum_{i=0}^{L} \sum_{k, k'} \frac{1}{L} e^{i(k-k')i} e^{-ik'} \tilde{c}_k^\dagger \tilde{c}_{k'} = \sum_{k} e^{-ik} \tilde{c}_k^\dagger \tilde{c}_k.
\end{align}

Hence, the Hamiltonian can be written as
\begin{eqnarray}
  H = \sum_{k} (\epsilon_0 - 2t\cos k) \tilde{c}_k^\dagger \tilde{c}_k.
\end{eqnarray}

This gives the dispersion relation 
\begin{eqnarray}
  E_k = \epsilon_0 - 2t\cos k.
\end{eqnarray}
and coefficients
\begin{eqnarray}
  \alpha_{kj} = \frac{1}{\sqrt{L}} e^{-ikj}.
\end{eqnarray}
The $k\in \frac{2\pi}{L}\mathbb{Z}/L\mathbb{Z}$ means the lattice momentum of the particle created by $\tilde{c}_k^\dagger$.

For the limit $N_e \ll L$, The dispersion relation on the region including Fermi points $k_F$ and $-k_F$ can be approximated by quadratic dispersion relation:
\begin{eqnarray}
  E_k = \epsilon_0 - 2t\cos k \approx \epsilon_0 - 2t + t k^2.
\end{eqnarray}

\subsection{(2b)}

The ground state of the fermionic many-body system is given by filling the single-particle states with lowest energy. For spin-1/2 case, there are two degenerate states for each momentum $k$. Hence, for a five-particle system, there are four degenerate ground states, which are formed by placing two particles at $k=0$ and distributing the remaining three particles among the four available states with $k=\pm\frac{2\pi}{L}$ and $\sigma=\uparrow,\downarrow$. Their ground state energy is given by
\begin{equation}
  E_0 = 5\epsilon_0 - 4t - 6t\cos\frac{2\pi}{L} \approx 5\epsilon_0 - 10t + \frac{12\pi^2 t}{L^2}.
\end{equation}

And the first excited state is $4\times 6 = 24$-fold degenerate, which is formed by placing two particles at $k=0$, two particles at $k=\pm\frac{2\pi}{L}$, and the remaining particle at $k=\pm\frac{4\pi}{L}$ with $\sigma=\uparrow,\downarrow$. The energy of the first excited state is given by
\begin{equation}
  E_1 = 5\epsilon_0 - 4t - 4t\cos\frac{2\pi}{L} - 2t\cos\frac{4\pi}{L} \approx 5\epsilon_0 - 10t + \frac{24\pi^2 t}{L^2}.
\end{equation}

\end{document}
